\chapter{Covers and Connection Descent}
\label{ch:neighborhood-parents}

This chapter asks how the pointwise parent construction of the main manuscript,
whose datum at one context is \cref{thm:rg-pointwise-parent-datum} and whose free
energy obeys \cref{thm:rg-exact-coarse-vfe}, extends over a neighborhood of
contexts, and when fine sections, laws, and connections descend along a base map.
The substantive answer is the descent criterion with its two obstructions
(\cref{thm:rg-cover-descent-obstructions}): a section descends exactly when its
image is constant on the fibers of the base map, and a connection form exactly
when its two pullbacks to the fiber product agree. Compatibility of local laws,
gluing of availability parents, and descent of pulled-back data hold by
construction once one ambient law is adopted (\cref{thm:rg-cover-compatible-laws}).
The geometric results show that a coarse map with context-dependent weights
keeps a first-order connection defect (\cref{thm:rg-cover-nonflat-defect}), that
a metric projection can create curvature from a flat source and that staged
projections compose only when the intermediate quotient metric is propagated
(\cref{thm:rg-cover-projected-connection,thm:rg-cover-quotient-composition}),
and that a normalized record distinguishes two base holonomy classes
(\cref{thm:rg-cover-operational-holonomy}). That record is the base-loop
companion of the agent-network witness \cref{thm:cg-loop-operational-holonomy}
of the main manuscript, which its interpretation chapter uses in
\cref{open:phil-testability}; the covariant defect used throughout is that of
\cref{def:geo-covariant-defects}.

The chapter declines to claim that the ambient law is inferred from independently
specified local marginals; that the full fine Gaussian family forms one
fixed-rank statistical manifold across a change of active set, or that a parent
extends through a context at which all its children are unavailable; that a
comparison across coarse levels gives Fisher contraction for derivatives in the
context, an inference-time decrease, or increased evidence, since the channels
depend on the context and differentiating in \(c\) also differentiates the
channel; that any construction here recovers a connection from data, identifies
an arbitrary graph-link product with base transport, or infers nontrivial bundle
topology from connection holonomy, which a connection on a trivial bundle can
already carry; that the change formulas of \cref{prop:pb-pullback-connection-change}
or the passive covariance of \cref{thm:pb-pullback-gauge-invariance} select a
canonical connection; and that the holonomy-sensitive record is a canonical
observable or has an empirical realization. \status{NOT-CLAIMED} The open
constructions are the gluing of independently supplied local data
(\cref{open:rg-cover-local-gluing}), quotients beyond the chosen metric
splitting (\cref{open:rg-cover-quotient-extension}), informative base
aggregation (\cref{open:rg-cover-base-aggregation}), graph-to-base identification
(\cref{open:rg-cover-graph-base-identification}), and connection recovery and
record-class sensitivity (\cref{open:rg-cover-holonomy-recovery}).

\section{A cover without a common intersection}
\label{sec:rg-cover-domains}

\definitionheading{Four-agent cover and availability weights}{def:rg-cover-four-agent}
Take the contextual base \(\mathcal U=(-3/2,3/2)^2\) with coordinates
\(c=(c^1,c^2)\) and four agents labeled \(0,1,2,3\), a local indexing used only
in this example. Their domains are
\begin{equation}
\begin{aligned}
 \mathscr I_0&=(-8/5,-1/10),& \mathscr I_1&=(-8/5,4/5),\\
 \mathscr I_2&=(-4/5,8/5),& \mathscr I_3&=(1/10,8/5),\\
 \mathcal U_i&=\mathcal U\cap(\mathscr I_i\times\R).
\end{aligned}
\label{eq:rg-cover-open-domains}
\end{equation}
No probability measure is placed on the base. With \(\mathrm m_i\) and
\(\mathrm h_i\) the midpoint and half-width of \(\mathscr I_i\), the availability
weights are
\begin{equation}
 w_i(c)=
 \begin{cases}
 \exp(1-1/t_i(c)),&t_i(c)>0,\\
 0,&t_i(c)\leq0,
 \end{cases}
 \qquad t_i(c)=1-\left(\frac{c^1-\mathrm m_i}{\mathrm h_i}\right)^2.
\label{eq:rg-cover-bumps}
\end{equation}
They encode availability, not posterior source probabilities: \(w_i(c)\) is a
per-endpoint retention probability, applied conditionally independently to every
record endpoint of agent \(i\), so one agent can be retained in one record and
replaced by the null endpoint \(\bot\) in another at the same context. A single
presence variable per agent, applied to all of its endpoints, would be a
different mechanism.

\theoremheading{The overlap pattern}{thm:rg-cover-nerve}
The sets \(\mathcal U_i\) cover \(\mathcal U\). Their nonempty pairwise
intersections have indices \(01,02,12,13,23\), their nonempty triple
intersections have indices \(012,123\), and their fourfold intersection is empty.
In particular,
\begin{equation}
 \mathcal U_{012}=\mathcal U\cap\{-4/5<c^1<-1/10\},\qquad
 \mathcal U_{123}=\mathcal U\cap\{1/10<c^1<4/5\}.
\label{eq:rg-cover-triples}
\end{equation}
The cover has two nonempty triple overlaps and no common intersection.
\status{ESTABLISHED}

\noindent\emph{Proof.}
An intersection of the intervals is nonempty exactly when its largest lower
endpoint is smaller than its smallest upper endpoint, which gives the listed
intersections. The intervals \(\mathscr I_1,\mathscr I_2\) alone cover the first
coordinate of \(\mathcal U\), and \(\mathcal U_0\cap\mathcal U_3=\varnothing\)
excludes the fourfold intersection. \(\square\)

\theoremheading{Smooth availability with an explicit null event}{thm:rg-cover-availability}
Each \(w_i\) is smooth, lies in \([0,1]\), and is positive exactly on
\(\mathcal U_i\). Adjoin the null endpoint \(\bot\) and define the endpoint channel
\begin{equation}
 \mathsf K^{\mathrm{av}}_c(i,j)=w_i(c)\mathbf1_{j=i}+(1-w_i(c))\mathbf1_{j=\bot},
 \qquad \mathsf K^{\mathrm{av}}_c(\bot,j)=\mathbf1_{j=\bot}.
\label{eq:rg-cover-availability-channel}
\end{equation}
For any fixed finite number of directed records, conditionally independent
applications of this channel to their endpoints give a smooth normalized channel
on the complete record state, so smooth ambient event probabilities give smooth
visible joint event probabilities, including at agent entry and exit.
\status{ESTABLISHED}

\noindent\emph{Proof.}
On its positive domain every derivative of \(\exp(-1/t)\) is that exponential
times a polynomial in \(1/t\), and all such products tend to zero as
\(t\downarrow0\), so extension by zero is smooth; since \(0<t_i\leq1\) inside the
interval, the range follows. Each row of \(\mathsf K^{\mathrm{av}}_c\) sums to one,
products of endpoint rows remain normalized, and their entries and finite sums
against smooth ambient probabilities are smooth. The channel preserves any
dependence it does not remove; it does not replace correlated records by the
product of their marginals. \(\square\)

\propositionheading{Occupancy and conditional attention on the cover}{prop:rg-cover-zero-occupancy}
For each channel \(x\in\{b,m\}\) let \(\eta^x_{ij}(c)\) be the marginal law of one
visible directed record, including \(\bot\), and set
\begin{equation}
 \alpha_i^x(c)=\sum_j\eta^x_{ij}(c),\qquad
 \beta^x_{ij}(c)=\frac{\eta^x_{ij}(c)}{\alpha_i^x(c)}
 \quad\text{only when }\alpha_i^x(c)>0.
\label{eq:rg-cover-conditional-attention}
\end{equation}
If every ambient receiver has positive mass, then \(\alpha_i^x>0\) exactly on
\(\mathcal U_i\), because the visible occupancy of receiver \(i\) equals its
ambient receiver mass times \(w_i\), and the conditional row \(\beta^x_{i\cdot}\)
is smooth there. At zero occupancy every entry of the event-law row of \(i\) is
zero, so the event law imposes no conditional row, whatever conditional
completion is chosen; smoothness of \(\eta^x\) and \(\alpha^x\) therefore
supplies no identifiable attention at zero occupancy. \status{ESTABLISHED}

\noindent\emph{Proof.}
The receiver endpoint of a record survives with probability \(w_i(c)\),
independently of the rest of the record, so summing the visible record law over
sources, \(\bot\) included, gives the ambient receiver mass times \(w_i(c)\),
which is positive exactly where \(w_i\) is. A quotient of smooth functions with
positive denominator is smooth. If \(\alpha_i^x(c)=0\), each nonnegative summand
\(\eta^x_{ij}(c)\) vanishes. No null outcome is removed or renormalized in this
calculation. \(\square\)

The joint event law, the occupancy, and the conditional row remain different
objects, as at one context; a generative posterior row is obtained from its own
event law by the same operation.

\section{Compatible laws, availability parents, and pullback descent}
\label{sec:rg-cover-law-descent}

\hypothesisheading{Ambient law on the cover}{hyp:rg-cover-ambient-law}
For the compatibility statements of this chapter, one smooth normalized ambient
joint
\(\mathbb P_c(dY,do,de\mid X)\), a smooth selected recognition
\(\mathbb Q_{c,o,X}(dY,de)\), and a reference law are given on one finite product
of agent fibers and record states for every \(c\in\mathcal U\), with fixed
structural datum \(X\), positive finite evidence at the admitted observation, and
selected posterior versions pushed with these same laws. The continuous
coordinates may be correlated Gaussian coordinates with positive-definite ambient
covariance, and records have finite laws; the section of agent \(i\) is exposed
only on \(\mathcal U_i\). This excludes local laws supplied on the members of the
cover without a common ambient extension; those are the subject of
\cref{open:rg-cover-local-gluing}.

\definitionheading{Availability parent}{def:rg-cover-availability-parent}
For agent subsets \(V_1\subseteq V_2\) let \(\mathrm{pr}_{V_1\leftarrow V_2}\)
project continuous coordinates to \(V_1\) and send every record endpoint outside
\(V_1\) to \(\bot\), with null endpoints kept null; restriction of the base
variable to a smaller open set is a separate operation. For a block \(I_A\) the
fixed-child domain, the availability domain, and the availability mass are
\begin{equation}
 \mathcal U_A^{\cap}=\bigcap_{i\in I_A}\mathcal U_i,\qquad
 \mathcal U_A^{\cup}=\bigcup_{i\in I_A}\mathcal U_i,
 \qquad s_A(c)=\sum_{i\in I_A}w_i(c).
\label{eq:rg-cover-parent-domains}
\end{equation}
The pointwise parent of the main manuscript lives on \(\mathcal U_A^{\cap}\). On
\(\mathcal U_A^{\cup}\) the availability parent uses the convex coefficients
\(\widehat w_{Ai}=w_i/s_A\) and all locally available joint child data supplied by
the ambient law; these coefficients are a separate choice, not derived from the
thinned record law of \eqref{eq:rg-cover-availability-channel}. For the explicit
geometric example take the standard two-dimensional representation of
\(\mathrm{SO}(2)\) separately for belief and model, let \(O_{\mathrm f},O_A\) be the
resulting orthogonal frame matrices of the active fine and parent coordinates,
and set
\begin{equation}
 L_A=(\widehat w_{Ai}I_4)_{i\in I_A},\qquad
 F_A=O_AL_AO_{\mathrm f}^{-1},
\label{eq:rg-cover-parent-channel}
\end{equation}
a deterministic normalized channel on the continuous coordinates whose event
channel sends a child endpoint to its block label and keeps \(\bot\). The example
uses equal standard representations; an inequivalent pair needs its own
intertwiners.

\definitionheading{Pullback presentation along a base submersion}{def:rg-cover-pullback-presentation}
Let \(f:\widetilde{\mathcal U}\to\mathcal U\) be a smooth surjective submersion,
let the principal and associated bundles over both bases be trivial, and let
\(O:\widetilde{\mathcal U}\to\mathrm{SO}(2)\) be a smooth passive frame field,
with \(\mathbf O\) repeating its representation in each separately typed
associated coordinate. The principal and associated maps
\begin{equation}
 \Psi_{\mathscr P}\bigl(\tilde c,u\bigr)=\bigl(f(\tilde c),O(\tilde c)^{-1}u\bigr),\qquad
 \Psi\bigl(\tilde c,v\bigr)=\bigl(f(\tilde c),\mathbf O(\tilde c)^{-1}v\bigr)
\label{eq:rg-cover-base-maps}
\end{equation}
cover \(f\). Given target sections, laws, connection potentials \(A_{\mathrm t}\),
and channels \(F_{\mathrm t}\) on \(\mathcal U\), their source presentations are
defined by \(\mathbf O\): sections and laws are pulled back by \(f\) and
re-expressed by \(\mathbf O\), and
\begin{equation}
 A_{\mathrm s}=\mathbf Of^*A_{\mathrm t}\mathbf O^{-1}-d\mathbf O\mathbf O^{-1},\qquad
 F_{\mathrm s}=\mathbf O_{\mathrm c}f^*F_{\mathrm t}\mathbf O_{\mathrm f}^{-1},
\label{eq:rg-cover-base-pullback}
\end{equation}
with observations and structural data pulled back unchanged.

\propositionheading{A global law induces compatible local data, and pullbacks descend}{thm:rg-cover-compatible-laws}
Assume \cref{hyp:rg-cover-ambient-law}.
\begin{enumerate}
\item[(a)] For \(V_1\subseteq V_2\subseteq V_3\) and on nested open context domains,
\begin{equation}
 \mathrm{pr}_{V_1\leftarrow V_2}\mathrm{pr}_{V_2\leftarrow V_3}=\mathrm{pr}_{V_1\leftarrow V_3},
 \qquad
 (\mathbb Q_{V_3}\mathrm{pr}_{V_2\leftarrow V_3})\mathrm{pr}_{V_1\leftarrow V_2}
 =\mathbb Q_{V_3}\mathrm{pr}_{V_1\leftarrow V_3},
\label{eq:rg-cover-restriction-law}
\end{equation}
and the same holds for the generative joint, the reference, and the selected
posterior families; smooth local presentations that obey their transition maps
on pairwise overlaps obey the corresponding triple-overlap cocycles.
\item[(b)]\label{thm:rg-cover-parent-gluing} The availability parent of
\cref{def:rg-cover-availability-parent} is smooth on \(\mathcal U_A^{\cup}\). For
disjoint blocks the active rows of the combined channel are independent, so a
positive-definite fine Gaussian covariance pushes to a positive-definite
covariance on the active parent fiber. If a second stage weights the active
blocks by \(s_A/\sum_Bs_B\), the direct fine coefficient is
\begin{equation}
 \frac{s_A}{\sum_Bs_B}\frac{w_i}{s_A}
 =\frac{w_i}{\sum_jw_j}\quad(i\in I_A,\ s_A>0),
\label{eq:rg-cover-staged-weights}
\end{equation}
so the direct and staged complete law channels agree. For the blocks \(01\) and
\(23\) the final availability parent exists on all of \(\mathcal U\), whereas a
fixed four-child parent has empty domain.
\item[(c)]\label{thm:rg-cover-base-descent} For a pullback presentation
(\cref{def:rg-cover-pullback-presentation}), sections, full laws, evaluators, and
channels descend through \(\Psi\). Both channel connections are compatible with
descent, \(d\Psi+f^*A_{\mathrm t}\Psi-\Psi A_{\mathrm s}=0\), their curvatures
are \(\operatorname{Curv}(A_{\mathrm s})=\mathbf Of^*\operatorname{Curv}(A_{\mathrm t})\mathbf O^{-1}\)
with \(\operatorname{Curv}(A)=dA+A\wedge A\), and the vertical coarse defect of
\(F_{\mathrm s}\) is the conjugate pullback of that of \(F_{\mathrm t}\). The cover
identities persist on the preimages \(f^{-1}(\mathcal U_i)\).
\end{enumerate}
Each clause holds by construction, from functoriality of pushforward, the frame
cocycle, and the pullback definitions. \status{ESTABLISHED}

\noindent\emph{Proof.}
(a) Coordinate projections compose, and sending a removed endpoint to \(\bot\)
while keeping \(\bot\) also composes, so pushforward through these deterministic
channels gives \eqref{eq:rg-cover-restriction-law}; restricting in \(c\) twice is
restriction to the smaller domain. For three local frame matrices
\(O_1,O_2,O_3\) of one associated fiber the coordinate transformations are
\(O_kO_j^{-1}\), and \((O_3O_2^{-1})(O_2O_1^{-1})=O_3O_1^{-1}\); pushing all
coordinates and full covariances by these transformations gives the law cocycle.
These are presentations of one law, not equalities between beliefs of different
agents.

(b) On \(\mathcal U_A^{\cup}\) we have \(s_A>0\), so every coefficient is smooth;
at a child boundary its vanishing bump extends smoothly while another child keeps
the denominator positive, and the formula is the same in every chart. Each active
block row has nonzero coefficients on its own disjoint child block, which gives
full row rank, and for a nonzero parent vector \(v\),
\(v^\top F\Sigma F^\top v=(F^\top v)^\top\Sigma(F^\top v)>0\). At two stages the
coefficient product is \eqref{eq:rg-cover-staged-weights}; an absent block
contributes zero directly and needs no conditional coefficient, frame matrices
cancel at intermediate fibers, endpoint maps compose with the same null
convention, and equal complete channels push every law equally. By
\cref{thm:rg-cover-nerve} \(\mathcal U_0\cap\mathcal U_3=\varnothing\), while
\(\mathcal U_0\cup\mathcal U_1=\{c^1<4/5\}\cap\mathcal U\) and
\(\mathcal U_2\cup\mathcal U_3=\{c^1>-4/5\}\cap\mathcal U\) cover \(\mathcal U\).

(c) The principal map is equivariant, since
\(O^{-1}(uk)=(O^{-1}u)k\). Pulling back the target principal connection form
\(u^{-1}A_{\mathrm t}u+u^{-1}du\) by \(\Psi_{\mathscr P}\) gives
\eqref{eq:rg-cover-base-pullback}, and differentiating \(\Psi=\mathbf O^{-1}\)
gives \(d\Psi=-\mathbf O^{-1}d\mathbf O\mathbf O^{-1}\), which cancels the
remaining terms of \(f^*A_{\mathrm t}\Psi-\Psi A_{\mathrm s}\). Expanding
\(dA_{\mathrm s}+A_{\mathrm s}\wedge A_{\mathrm s}\), the derivatives of
\(\mathbf O\) cancel and leave the conjugate pullback curvature; the same
expansion of \(dF_{\mathrm s}+A_{\mathrm s,\mathrm c}F_{\mathrm s}-F_{\mathrm s}A_{\mathrm s,\mathrm f}\)
leaves the conjugate pullback of the target defect. A section
\(s_{\mathrm s}=\mathbf O f^*s_{\mathrm t}\) satisfies
\(\Psi\circ s_{\mathrm s}=s_{\mathrm t}\circ f\), which is constant on the fibers
of \(f\). For a target kernel \(C_c\), the source kernel undoes its input frame,
applies \(C_c\), and restores its output frame; evaluating on a target measurable
set gives
\begin{equation}
 \Psi_{\mathrm c\#}(\mathbb Q_{\mathrm s}C_{\mathrm s})
 =(\Psi_{\mathrm f\#}\mathbb Q_{\mathrm s})C_{\mathrm t}.
\label{eq:rg-cover-base-channel-square}
\end{equation}
The same change of variables holds for the full generative joint, the reference,
and the selected posterior versions. Keeping observation coordinates unchanged
preserves evidence, common invertible presentation maps preserve relative entropy
because the density Jacobians cancel, and the normalized evaluator transforms by
its separately typed input-output square. Restrictions to \(\mathcal U_i\) pull
back to restrictions to \(f^{-1}(\mathcal U_i)\). \(\square\)

At each fixed admitted context, every normalized recognition-independent state
channel of this section is covered by \cref{thm:rg-exact-coarse-vfe}.

\openproblemheading{Gluing of independently supplied local data}{open:rg-cover-local-gluing}
\Cref{thm:rg-cover-compatible-laws} starts from one ambient law and therefore
cannot establish the converse. What is not settled is which local laws, channels,
and presentations supplied on the members of a cover without an ambient law admit
one joint extension compatible with every restriction and selected conditional
version. The same question is open for the descent of independently specified
source data along an informative base map and for a canonical choice of ambient
dependence, channel, or membership. The problem is settled by a gluing theorem
with hypotheses on the local data alone, or by an obstruction computed from them;
a compatible ambient law supplies one extension but does not reconstruct every
marginal family.

\section{Nonflat connections and metric-dependent projections}
\label{sec:rg-cover-connections}

\definitionheading{Reference potentials}{def:rg-cover-nonflat-potentials}
On the trivial common principal \(\mathrm{SO}(2)\) bundle over \(\mathcal U\),
choose two connections independently, with local reference potentials
\begin{equation}
 A^x_*=\frac{\kappa_x}{2}\bigl(c^1dc^2-c^2dc^1\bigr)\mathbb J,\qquad
 \mathbb J=\begin{pmatrix}0&-1\\1&0\end{pmatrix},\qquad x\in\{b,m\}.
\label{eq:rg-cover-nonflat-potential}
\end{equation}
The convention is \(\nabla=d+A\), so the potential in a frame \(O\) is
\(A=OA_*O^{-1}-dOO^{-1}\); this is the represented form of
\eqref{eq:geo-local-connection-b} and \eqref{eq:geo-local-connection-m}, with
\(O\) the represented inverse of the passive reframing of the channel. The
connections are supplied geometric data, not quantities determined by attention.

\theoremheading{A nonflat parent retains its connection defect}{thm:rg-cover-nonflat-defect}
The reference curvatures are
\(\operatorname{Curv}(A^x_*)=\kappa_x\mathbb Jdc^1\wedge dc^2\). Wherever the weights
in \(L\) are smooth, in particular on \(\mathcal U_A^{\cup}\), and with the same
channel connection repeated in its child fibers, the channel
\(F=O_{\mathrm c}LO_{\mathrm f}^{-1}\) has covariant defect
\begin{equation}
 \mathcal DF=dF+A_{\mathrm c}F-FA_{\mathrm f}
   =O_{\mathrm c}(dL)O_{\mathrm f}^{-1},
\label{eq:rg-cover-connection-defect}
\end{equation}
and for composable maps
\begin{equation}
 \mathcal D(F_2F_1)=(\mathcal DF_2)F_1+F_2(\mathcal DF_1).
\label{eq:rg-cover-defect-composition}
\end{equation}
Nonzero curvature is therefore compatible with gauge covariance, and a coarse map
with context-dependent weights is not parallel. \status{ESTABLISHED}

\noindent\emph{Proof.}
The abelian potential has \(A_*\wedge A_*=0\), and differentiating its
coefficient gives \(\kappa_x\mathbb Jdc^1\wedge dc^2\). Differentiate
\(F=O_{\mathrm c}LO_{\mathrm f}^{-1}\) and substitute the frame potentials; the
frame derivatives cancel, and the repeated reference connection commutes through
the scalar block weights, leaving \(dL\). The product rule gives
\eqref{eq:rg-cover-defect-composition} after adding and subtracting
\(F_2A_1F_1\). \(\square\)

\definitionheading{Operational base-holonomy witness}{def:rg-cover-holonomy-witness}
An operational base-holonomy witness consists of a named principal bundle and a
family of connections, an assigned base loop, a gauge-invariant statistic of a
normalized population record, a typed map from the connection data to the record
law, and two admissible connection data for which every input other than the
connection is fixed, the represented loop holonomies have distinct conjugacy
classes, and the induced laws of the statistic differ. It is strong when the
record-law map factors nontrivially through the conjugacy class of the
represented loop holonomy.

\theoremheading{An operational record sensitive to a declared base holonomy}{thm:rg-cover-operational-holonomy}
Fix the base point \(c_\circ=(7/20,1/5)\) and the closed rectangular path with
successive vertices \((7/20,1/5)\), \((13/20,1/5)\), \((13/20,1/2)\),
\((7/20,1/2)\), and \(c_\circ\), which lies in the domain \(\mathcal U_2\) of the
probe agent; only its fibers at \(c_\circ\) enter the record. Choose a selector
\(P_{\rm pr}\ne0\) onto its separately typed belief and model coordinates and a
positive-definite noise covariance \(R_{\rm pr}\). Let a finite label \(\iota\)
prescribe connection parameters \((\kappa_{b,\iota},\kappa_{m,\iota})\) in
\eqref{eq:rg-cover-nonflat-potential}, with positive prior masses
\(\mathsf p_\iota\). Keep one normalized joint \(p_0(Y,e,o_0)\), independent of
\(\iota\), with matter posterior \(\boldsymbol\Pi_0(dY\mid o_0)\) at an admitted
\(o_0\), let \(\mathbf O_\circ\) be the represented probe frame at \(c_\circ\) in
both channels, and declare the additional record \(o_1\) and joint by
\begin{equation}
\begin{aligned}
 \operatorname{Hol}_\iota&=\mathbf O_\circ\exp\left[-\mathcal A
       \operatorname{diag}(\kappa_{b,\iota}\mathbb J,\kappa_{m,\iota}\mathbb J)\right]\mathbf O_\circ^{-1},
 \qquad \mathcal A=9/100,\\
 B_\iota&=(\operatorname{Hol}_\iota-I_4)P_{\rm pr},\\
 p(\iota,Y,e,o_0,o_1)&=\mathsf p_\iota p_0(Y,e,o_0)
                       \mathcal N(o_1;B_\iota Y,R_{\rm pr}).
\end{aligned}
\label{eq:rg-cover-loop-observation}
\end{equation}
This is a normalized statistical experiment, and given \(\iota\) and \(o_0\) the
record is distributed as \(B_\iota Y+\varepsilon\), with
\(Y\sim\boldsymbol\Pi_0(\cdot\mid o_0)\) and independent
\(\varepsilon\sim\mathcal N(0,R_{\rm pr})\). The record laws of two labels
coincide exactly when \(B_\iota Y\) and \(B_{\iota'}Y\) have the same law under
\(\boldsymbol\Pi_0(\cdot\mid o_0)\). For the parameter pairs \((0,0)\) and
\((7/10,-2/5)\), labeled \(\iota=0,1\), the two record laws differ exactly when
the posterior law of \(P_{\rm pr}Y\) is not the point mass at \(0\), in particular
whenever \(\boldsymbol\Pi_0(\cdot\mid o_0)\) is Gaussian with positive-definite
covariance. Under the same condition the frame-invariant statistic
\(\mathsf S(o_1)=o_1^\top R_{\rm pr}^{-1}o_1\) has different laws under the two
labels, with means \(4\) and
\(4+\E_{\boldsymbol\Pi_0}\bigl[(B_1Y)^\top R_{\rm pr}^{-1}B_1Y\bigr]\in(4,\infty]\),
while the matter law, loop, selector, and noise are held fixed. Under that
condition these data form an operational base-holonomy witness in the sense of
\cref{def:rg-cover-holonomy-witness}. \status{ESTABLISHED}

\noindent\emph{Proof.}
The rectangle and the strip \(\mathscr I_2\times\R\) are convex, so every segment
stays in \(\mathcal U_2\). Integrating the reference one-form along a segment
from \(c_j\) to \(c_{j+1}\) gives \(\kappa\det(c_j,c_{j+1})/2\); summing the
segments gives \(\kappa\mathcal A\) with the signed area \(\mathcal A=9/100\), and
the transport equation \(\dot v=-A(\dot\gamma)v\) gives the displayed exponential,
which the endpoint frame conjugates at the base point. Each Gaussian factor
integrates to one in \(o_1\), and \(p_0\) and \(\mathsf p\) normalize the joint.

The supplied pairs give \(\operatorname{Hol}_0=I_4\) and \(B_0=0\), while
\(\operatorname{Hol}_1\) is conjugate to two plane rotations by the angles
\(-63/1000\) and \(9/250\); neither is a multiple of \(2\pi\), so
\(\operatorname{Hol}_1-I_4\) is invertible and the two holonomies have distinct
conjugacy classes. Under the declared joint the label is independent of
\((Y,e,o_0)\), so the joint is the record joint of
\cref{thm:cg-loop-operational-holonomy} with latent \((Y,e)\), observation
\(o_0\), prior masses \(\mathsf p_\iota\), probe coordinates \(P_{\rm pr}Y\in\R^4\),
represented holonomies \(\operatorname{Hol}_\iota\), and noise covariance
\(R_{\rm pr}\), and \(\mathsf S\) is its statistic \(\mathsf q\). The proofs of
parts (i) and (ii) there use the holonomies only as fixed matrices and apply
verbatim. They give the record kernel, the criterion for equal record laws, and
that the two record laws differ exactly when the posterior law of
\(P_{\rm pr}Y\) is not the point mass at \(0\), in which case the laws of
\(\mathsf S\) differ with the displayed means; otherwise both record laws, and
hence both laws of \(\mathsf S\), coincide. A Gaussian posterior with
positive-definite covariance gives \(P_{\rm pr}Y\) a nondegenerate law, which is
not a point mass. Under label \(0\) we have \(o_1=\varepsilon\), so \(\mathsf S\)
has the chi-square law with four degrees of freedom. In general the
label likelihood is
\(\mathrm{lik}_\iota(o_1)=\int\mathcal N(o_1;B_\iota y,R_{\rm pr})\boldsymbol\Pi_0(dy\mid o_0)\).
If the matter posterior is \(\mathcal N(\mu_0,\Sigma_0)\) with \(\Sigma_0\)
positive definite, Gaussian integration gives
\begin{equation}
\begin{gathered}
 \mathrm{lik}_\iota(o_1)=\mathcal N\bigl(o_1;B_\iota\mu_0,R_{\rm pr}+B_\iota\Sigma_0B_\iota^\top\bigr),\qquad
 Z(o_1,o_0)=Z_0(o_0)\sum_\iota\mathsf p_\iota\mathrm{lik}_\iota(o_1),\\
 \mathsf p(\iota\mid o_1,o_0)=\frac{\mathsf p_\iota\mathrm{lik}_\iota(o_1)}{\sum_{\iota'}\mathsf p_{\iota'}\mathrm{lik}_{\iota'}(o_1)},
\end{gathered}
\label{eq:rg-cover-loop-evidence}
\end{equation}
and the mean of \(\mathsf S\) under label \(1\) is
\(4+\Tr(R_{\rm pr}^{-1}B_1\Sigma_0B_1^\top)+\mu_0^\top B_1^\top R_{\rm pr}^{-1}B_1\mu_0\).

Under a passive change of fine and probe coordinates by represented orthogonal
frame maps \(O_{\mathrm f},O_{\mathrm{pr}}\), send \(P_{\rm pr}\) to
\(O_{\mathrm{pr}}P_{\rm pr}O_{\mathrm f}^{-1}\), \(\operatorname{Hol}_\iota\) to
\(O_{\mathrm{pr}}\operatorname{Hol}_\iota O_{\mathrm{pr}}^{-1}\), \(o_1\) to
\(O_{\mathrm{pr}}o_1\), and \(R_{\rm pr}\) to
\(O_{\mathrm{pr}}R_{\rm pr}O_{\mathrm{pr}}^\top\). Then
\(B_\iota'=O_{\mathrm{pr}}B_\iota O_{\mathrm f}^{-1}\), so the whole observation
kernel is pushed coherently and its evidence is invariant, and \(\mathsf S\) is
unchanged since
\(o_1^\top O_{\mathrm{pr}}^\top(O_{\mathrm{pr}}R_{\rm pr}O_{\mathrm{pr}}^\top)^{-1}O_{\mathrm{pr}}o_1
=o_1^\top R_{\rm pr}^{-1}o_1\). A zero-area loop gives \(B_\iota=0\) for every
label and no label reweighting. \(\square\)

The proof uses the base loop only to compute \(\operatorname{Hol}_\iota\); the
distinguishability argument needs only a represented group element acting on the
probe fibers at one point and conjugated by the probe frame under a passive
change. At \(c_*\) the represented loop product of independent graph links
around a closed walk of the interaction graph plays that role, which is
\cref{thm:cg-loop-operational-holonomy} of the main manuscript, with no base
curve, connection, or neighborhood. A specified connection is thereby associated
with a normalized population record; reading that association as a statement
about physical geometry is an interpretation, not a result of this chapter.

Identifying a graph link \(\Theta_e\) with a base transport \(\Omega_{\gamma_e}\)
requires an embedded interaction complex, assigned piecewise-smooth base curves,
and equality with the path-ordered exponential of a specified connection.
\Cref{hyp:geo-graph-base-transport} is satisfiable by declaration, by setting each
represented link equal to a chosen base transport, but that is a definition, not
a theorem; the cover of this chapter contains no interaction graph and no graph
links, and its base witness uses one declared loop for one probe agent. Even an
exact assignment concerns connection holonomy, not topological nontriviality.

\openproblemheading{Graph-to-base identification}{open:rg-cover-graph-base-identification}
No realization of independently specified graph links as base transports is
known. The problem is to choose vertex contexts, an embedding, and oriented edge
curves for given graph data and a connection, and to prove that the independently
specified represented links equal the corresponding base transports as in
\eqref{eq:geo-link-pt-hypothesis}, together with hypotheses on \(\mathcal C\),
\(G\), and the curve class under which the realization is unique up to gauge,
or to exhibit graph data for which no such choice exists.

\openproblemheading{Connection recovery and record-class sensitivity}{open:rg-cover-holonomy-recovery}
\Cref{thm:rg-cover-operational-holonomy} supplies one witness for Gaussian record
noise and one pair of connection labels. What is not settled is whether a
connection, or its holonomy class, can be recovered from population records
beyond such a controlled family, and for which record classes holonomy
sensitivity persists. The problem is settled positively by a recovery theorem
for a declared record class, and negatively by a theorem that every statistic in
a declared admissible controlled class is insensitive to the holonomy class.

\definitionheading{Minimum-norm splitting and projected connection}{def:rg-cover-minimum-norm-splitting}
Let \(E\) carry a Euclidean fiber metric and a compatible connection \(d+A\), with
\(A^\top=-A\) in orthonormal frames, and let \(F:E\to E_{\mathrm c}\) be a smooth
surjective vector-bundle map of fixed rank. Set
\begin{equation}
\begin{gathered}
 \bar{\mathsf g}=(FF^\top)^{-1},\qquad F^\dagger=F^\top\bar{\mathsf g},\qquad
 P_F=F^\dagger F,\qquad P_F^\perp=I-P_F,\\
 A^{\mathrm{quot}}=F(dF^\dagger+AF^\dagger)=(FA-dF)F^\dagger.
\end{gathered}
\label{eq:rg-cover-quotient-data}
\end{equation}
Here \(P_F\) is a fiberwise orthogonal projector, \(\bar{\mathsf g}\) is the
quotient metric of the chosen metric, \(F^\dagger\) is the minimum-norm lift, and
\(d+A^{\mathrm{quot}}\) is the projected connection.

\theoremheading{Projected connection versus exact connection descent}{thm:rg-cover-projected-connection}
The projected connection preserves \(\bar{\mathsf g}\). With the second fundamental
form \(\mathrm{II}=P_F^\perp(dF^\dagger+AF^\dagger)\) and the defect
\(\mathcal DF=dF+A^{\mathrm{quot}}F-FA\),
\begin{equation}
 d\bar{\mathsf g}=(A^{\mathrm{quot}})^\top\bar{\mathsf g}+\bar{\mathsf g}A^{\mathrm{quot}},\qquad
 \mathcal DF=(dF-FA)P_F^\perp=\bar{\mathsf g}^{-1}\mathrm{II}^\top,
\label{eq:rg-cover-projected-defect}
\end{equation}
and \(F\) is parallel exactly when \(\ker F\) is preserved by the fine connection.
The curvature is
\begin{equation}
 \operatorname{Curv}(A^{\mathrm{quot}})=F\operatorname{Curv}(A)F^\dagger
 +\bar{\mathsf g}^{-1}\mathrm{II}^\top\wedge\mathrm{II},
 \qquad \operatorname{Curv}(A)=dA+A\wedge A,
\label{eq:rg-cover-projected-curvature}
\end{equation}
so a projected connection can be curved with a flat source: for
\begin{equation}
 A=0,\qquad F(c^1,c^2)=\begin{pmatrix}1&0&-c^1\\0&1&-c^2\end{pmatrix},
 \qquad \operatorname{Curv}(A^{\mathrm{quot}})(0)=-\mathbb Jdc^1\wedge dc^2\ne0.
\label{eq:rg-cover-projected-curvature-example}
\end{equation}
For the scalar-weight parent \(L=(\widehat w_iI)_i\) with a common internal
connection \(A_*\), the construction instead gives
\begin{equation}
 \bar{\mathsf g}=|\widehat w|^{-2}I,\qquad
 A^{\mathrm{quot}}=A_*-\tfrac12d\log|\widehat w|^2I,\qquad
 \operatorname{Curv}(A^{\mathrm{quot}})=\operatorname{Curv}(A_*),
 \qquad |\widehat w|^2=\sum_i\widehat w_i^2.
\label{eq:rg-cover-scalar-weight-quotient}
\end{equation}
\status{ESTABLISHED}

\noindent\emph{Proof.}
Direct multiplication gives \(FF^\dagger=I\), \(P_F=P_F^\top=P_F^2\), and
\((F^\dagger)^\top F^\dagger=\bar{\mathsf g}\). A fine representative of \(y\) is
uniquely \(F^\dagger y+v\) with \(v\in\ker F\), of squared norm
\(y^\top\bar{\mathsf g}y+\|v\|^2\), which identifies the quotient metric; the two
expressions for \(A^{\mathrm{quot}}\) agree by differentiating \(FF^\dagger=I\).
Multiplying by \(\bar{\mathsf g}\) gives
\(\bar{\mathsf g}A^{\mathrm{quot}}=(F^\dagger)^\top(dF^\dagger+AF^\dagger)\), and
adding its transpose cancels the skew \(A\) terms and gives \(d\bar{\mathsf g}\).
Differentiating \(FF^\dagger=I\) gives the first formula for \(\mathcal DF\). With
\(D_AP_F=dP_F+[A,P_F]\), differentiating \(P_F^2=P_F\) shows that \(D_AP_F\) has
only off-diagonal blocks, metric compatibility makes its coefficients symmetric,
so \((D_AP_F)F^\dagger=\mathrm{II}\) and
\(\mathcal DF=F(D_AP_F)P_F^\perp=\bar{\mathsf g}^{-1}\mathrm{II}^\top\). Its
vanishing is equivalent to preservation of the retained subbundle and, by
differentiating orthogonal inner products, of its orthogonal complement
\(\ker F\). On a retained section \(u=P_Fu\) the projected connection is
\(P_F\nabla\); expanding its square gives
\(P_F\operatorname{Curv}(A)u+P_F(D_AP_F)\wedge(D_AP_F)u\), since the term with
\(P_F(D_AP_F)P_F\) vanishes and \(P_F^\perp\nabla u=(D_AP_F)u\), and transporting
through \(F^\dagger\) gives \eqref{eq:rg-cover-projected-curvature}. The second
term is an endomorphism-valued two-form, not a nonnegative scalar energy. In the
example \(\bar{\mathsf g}(0)=I_2\) and the only nonzero row of \(\mathrm{II}(0)\) is
\((-dc^1,-dc^2)\), whose exterior square gives
\(-\mathbb Jdc^1\wedge dc^2\); all bundles there are trivial, so the curvature
does not indicate nontrivial topology. For the scalar-weight parent insert
\(F^\dagger=\operatorname{col}_i(\widehat w_iI)/|\widehat w|^2\): the additional
scalar one-form is exact, its square vanishes, and its mixed wedge terms with
\(A_*\) cancel. The quotient-orthonormal coordinates multiply \(y\) by
\(|\widehat w|^{-1}\) and remove this term; the variable-weight defect of
\cref{thm:rg-cover-nonflat-defect} need not vanish. \(\square\)

\theoremheading{Composition requires propagation of the quotient metric}{thm:rg-cover-quotient-composition}
For smooth surjections \(F_1:E_0\to E_1\) and \(F_2:E_1\to E_2\), start with the
metric \(I\) on \(E_0\), give \(E_1\) its induced metric
\(\bar{\mathsf g}_1=(F_1F_1^\top)^{-1}\), and use it in the second projection,
whose minimum-norm lift is
\(F_2^{\dagger_1}=\bar{\mathsf g}_1^{-1}F_2^\top(F_2\bar{\mathsf g}_1^{-1}F_2^\top)^{-1}\).
Then
\begin{equation}
 F_1^\dagger F_2^{\dagger_1}=(F_2F_1)^\top\bigl((F_2F_1)(F_2F_1)^\top\bigr)^{-1}=(F_2F_1)^\dagger,
\label{eq:rg-cover-quotient-lift-composition}
\end{equation}
and the quotient metrics and projected connections of the staged and direct
projections agree, whether or not either map is parallel. Resetting the
intermediate metric to \(I\) is a different construction and can change the
connection class: for \(A_0=0\),
\(F_1=\left(\begin{smallmatrix}2&0&0\\0&1&0\end{smallmatrix}\right)\), and
\(F_2=(1,c^1)\), the reset connection is \(-c^1dc^1/(1+(c^1)^2)\) and the direct
one is \(-c^1dc^1/(4+(c^1)^2)\), which differ by the exact form
\(d\bigl(\tfrac12\log[(4+(c^1)^2)/(1+(c^1)^2)]\bigr)\), so both are flat and
related by a rescaling of the quotient frame; for
\(F_1=\left(\begin{smallmatrix}2&0&0&1\\0&1&0&0\\0&0&3&0\end{smallmatrix}\right)\)
and \(F_2=\left(\begin{smallmatrix}1&c^1&0\\0&1&c^2\end{smallmatrix}\right)\), the
determinant of the \(dc^1\wedge dc^2\) coefficient of the curvature at
\((c^1,c^2)=(1/3,1/5)\) is \(15625/61731\) for the propagated metric and
\(3645/103823\) for the reset metric, so the reset connection is not gauge
equivalent to the direct one. \status{ESTABLISHED}

\noindent\emph{Proof.}
Multiplication gives \eqref{eq:rg-cover-quotient-lift-composition}, and the final
quotient metric has the same denominator. The product rule and
\(F_1F_1^\dagger=I\) give
\(F_2(d+A^{\mathrm{quot}}_1)F_2^{\dagger_1}=F_2F_1(d+A_0)(F_1^\dagger F_2^{\dagger_1})\),
which is the direct connection. The two examples are direct computations; gauge
equivalent connections have pointwise conjugate curvatures and hence equal
curvature determinants. \(\square\)

\openproblemheading{Quotients beyond the chosen metric splitting}{open:rg-cover-quotient-extension}
\Cref{thm:rg-cover-projected-connection,thm:rg-cover-quotient-composition} close
the chosen minimum-norm construction for a declared metric and fixed rank. Not
settled are a canonical effective-support quotient for arbitrary
positive-semidefinite edge weights, an intertwiner into an arbitrarily prescribed
associated representation, and quotient geometry across changes of rank. The
problem is settled by constant-rank support data, compatible represented actions,
and an admitted quotient or bundle map for which these objects are constructed,
or by an obstruction to them.

\section{Base descent with inherited nonzero curvature}
\label{sec:rg-cover-base-descent}

The pullback presentation of \cref{def:rg-cover-pullback-presentation} with
\(\widetilde{\mathcal U}=\mathcal U\times(-1,1)\), extra coordinate \(z\),
\(f(c,z)=c\), and \(O(z)=\exp(\mathrm kz\mathbb J)\) removes an explicitly
redundant base coordinate; it does not average contexts. By
\cref{thm:rg-cover-compatible-laws} its source data descend. The following
criterion concerns data that are given on the source rather than defined as
pullbacks.

\theoremheading{The descent criterion and its obstructions}{thm:rg-cover-descent-obstructions}
Let \(f:\mathcal M\to\overline{\mathcal M}\) be a smooth surjective submersion and
\(\Psi:E\to\overline E\) a smooth bundle map over \(f\). A smooth section \(s\) of
\(E\) has a unique descended section exactly when \(\Psi\circ s\) is constant on
every fiber of \(f\). For a smooth surjective principal submersion
\(\Psi_{\mathscr P}:\mathscr P\to\overline{\mathscr P}\), a connection form
descends exactly when its pullbacks by the two projections of
\(\mathscr P\times_{\overline{\mathscr P}}\mathscr P\) agree. In the example above
both criteria have concrete failures: the source law
\(\mathcal N(O(z)(1+z/4,0)^\top,I_2)\) is pushed by \(\Psi\) to the fiberwise law
\(\mathcal N((1+z/4,0)^\top,I_2)\), whose mean varies with \(z\), so it does not
descend, and no passive reframing makes it constant on the forgotten fiber; and
adding \(\mathrm sz\mathbb Jdc^1\), \(\mathrm s\ne0\), to the pullback potential
excludes compatible connection descent through this \(f\). By contrast the
pullback potentials \(A^x_{\mathrm s}=f^*A^x_*-\mathrm k\mathbb Jdz\) of
\eqref{eq:rg-cover-nonflat-potential} have zero mixed curvature and inherited
\(dc^1\wedge dc^2\) curvature \(\kappa_x\mathbb J\), so a nonzero coefficient in the
forgotten direction is compatible with exact descent. These are criteria for
given maps, not the existence of a canonical quotient. \status{ESTABLISHED}

\noindent\emph{Proof.}
Necessity follows by evaluating a descended object at two lifts of the same base
point or tangent vector. Conversely, a submersion has local smooth sections: in
local coordinates, augment its full-rank derivative by complementary source
coordinates and apply the inverse function theorem. Composing \(\Psi\circ s\) with
such a local section defines a smooth target section; fiber constancy makes the
definition independent of the local section and hence consistent on overlaps, and
surjectivity gives uniqueness. For forms, choose a local section \(\varsigma\) of
\(\Psi_{\mathscr P}\) and set \(\overline\omega=\varsigma^*\omega\); the
fiber-product equality applied to \((v,d\varsigma d\Psi_{\mathscr P}(v))\) gives
\(\omega=\Psi_{\mathscr P}^*\overline\omega\) and independence of \(\varsigma\).
Equivariance and reproduction of the principal vertical generators descend by
surjectivity and the equivariant group-fiber isomorphisms.

In the example the pushed mean is \(O(z)^{-1}O(z)(1+z/4,0)^\top\). Its relative
entropy to \(\mathcal N(0,I_2)\) is \((1+z/4)^2/2\), which every passive
reframing preserves, so no reframing makes the pair constant in \(z\). The
potential \(A^x_{\mathrm s}=f^*A^x_*-\mathrm k\mathbb Jdz\) follows from
\eqref{eq:rg-cover-base-pullback}; its \(dz\wedge dc^j\) curvature components
vanish because \(\mathbb J\) commutes with itself and the \(dz\) coefficient is
constant. Adding \(\mathrm sz\mathbb Jdc^1\) adds \(\mathrm s\mathbb Jdz\wedge dc^1\)
to the curvature, whose contraction with \(\partial_z\) is nonzero, whereas every
conjugate pullback curvature annihilates \(\partial_z\). \(\square\)

\openproblemheading{Informative base aggregation}{open:rg-cover-base-aggregation}
Choosing a probability measure on a varying base fiber and averaging its laws
defines a new mixture operation, which requires a base prior or reference, a joint
base-fiber model, and the induced observation weighting; neither fiber constancy
nor a canonical such measure follows from the cover. The problem is to construct
informative base aggregation of laws, channels, and connections along a base map
that collapses contexts, with a probability law on smooth random sections where
the construction needs one (compare \cref{open:prob-continuum-theory}), and to
state its descent and free-energy identities; it is settled by such a construction
or by an obstruction for a declared base map.
